\begin{center}
\textbf{\large{Abstract}}
\end{center}

\addcontentsline{toc}{chapter}{Abstract}
\begin{small}

Distilling a compact VLM from a heterogeneous teacher pool faces three
difficulties simultaneously: mismatched tokenizers make direct KD invalid,
teacher usefulness varies sample by sample and a routed teacher can still
produce gradients that oppose the supervised objective. Fixed teacher averaging
fails to address this third difficulty: it mixes helpful signals with
conflicting ones without discrimination.

\gracefull{} (\gracename{}) addresses all three. It combines Byte-Trie
Wasserstein knowledge distillation over a shared UTF-8 byte trie, sparse top-$k$
teacher routing conditioned on student hidden states and gradient-agreement
refinement of routed teacher weights. The Byte-Trie Wasserstein loss compares
teacher and student output distributions across tokenizers without shared-index
assumptions, while sparse routing and gradient-agreement refinement select
teachers per sample and down-weight those whose KD gradients conflict with
the CE objective.

Experiments on TextVQA, DocVQA and ChartQA evaluate \methodname{} with
SmolVLM-500M and SmolVLM-256M students with a four-teacher heterogeneous pool.
For SmolVLM-500M, \methodname{} reaches 69.38, 73.11 and 80.20 on TextVQA,
DocVQA and ChartQA, closing 55.03\%, 39.53\% and 55.34\% of the original
gap to the highest-scoring teacher. For SmolVLM-256M, \methodname{} reaches
60.73, 60.48 and 72.32 on the same benchmarks. In all comparisons,
\methodname{} improves over direct fine-tuning, same-family logit-based KD,
cross-family ULD and the heterogeneous multi-teacher baselines.


\vspace*{1cm}
\textbf{Keywords:} vision-language model, knowledge distillation,
multi-teacher distillation, sparse teacher routing, gradient agreement,
cross-tokenizer Wasserstein distillation
\end{small}
